% !PRINTED 325
\weberpage{325}
%<<CONT
heisst die Discriminante von $f$. Sie ist durch diese Definition nur bis auf einen numerischen Factor bestimmt, und kann bis jetzt nur in den einfachsten Fällen berechnet werden.

Die Punkte der Curve $f$, deren Coordinaten den Bedingungen (1) genügen, heissen \gesperrt{singuläre Punkte}.

Man kann auch die Frage aufwerfen, unter welcher Bedingung es möglich ist, dass auf einer Curve unendlich viele singuläre Punkte vorkommen. Damit dies eintritt, muss zunächst das Eliminationsresultat von einer der Unbekannten, etwa von $x_3$, aus den zwei Gleichungen $f'(x_1) = 0$, $f'(x_2) = 0$ identisch verschwinden, d.\,h. die beiden Functionen $f'(x_1)$ und $f'(x_2)$ müssen, als Functionen von $x_3$ betrachtet, einen gemeinschaftlichen Theiler haben, und nach Bd. I, \S.\,51 müssen sie daher auch einen (nicht constanten) gemeinsamen Theiler im Gebiete der Formen $x_1, x_2, x_3$ haben. Dieser Theiler muss dann wieder, wie aus denselben Erwägungen hervorgeht, einen gemeinschaftlichen Theiler mit $f$, oder, was dasselbe ist, mit $f'(x_3)$ haben, und es ergiebt sich also, dass nur dann unendlich viele singuläre Punkte vorhanden sein können, wenn die vier Formen $f, f'(x_1), f'(x_2), f'(x_3)$ einen gemeinschaftlichen Theiler haben.

Es sei nun $u$ ein solcher gemeinschaftlicher Theiler, den wir als \gesperrt{irreducibel} voraussetzen.

Wir setzen
\[
f(x_1, x_2, x_3) = uv
\]
\[
f'(x_1) = u\,\frac{\partial v}{\partial x_1} + v\,\frac{\partial u}{\partial x_1}
\]
\[
f'(x_2) = u\,\frac{\partial v}{\partial x_2} + v\,\frac{\partial u}{\partial x_2}
\]
\[
f'(x_3) = u\,\frac{\partial v}{\partial x_3} + v\,\frac{\partial u}{\partial x_3},
\]
und diese Gleichungen zeigen, da $u$ irreducibel ist und daher
\[
u, \quad \frac{\partial u}{\partial x_1}, \quad \frac{\partial u}{\partial x_2}, \quad \frac{\partial u}{\partial x_3}
\]
keinen gemeinschaftlichen Theiler haben können, dass $v$ durch $u$ theilbar sein muss; d.\,h. es können nur dann unendlich viele singuläre Punkte vorkommen, wenn $f(x_1, x_2, x_3)$ einen \gesperrt{quadratischen Theiler} hat. In diesem Falle sagen wir, dass die Curve $f$ einen doppelt zählenden Bestandtheil enthält.
